\documentclass{article}
\usepackage[portuguese]{babel}
\usepackage{listings}
\usepackage{xcolor}

\definecolor{codegreen}{rgb}{0,0.6,0}
\definecolor{codegray}{rgb}{0.5,0.5,0.5}
\definecolor{codepurple}{rgb}{0.8,0,0.2}
\definecolor{backcolour}{rgb}{.8,.8,1}

\lstdefinestyle{mystyle}{
    backgroundcolor=\color{backcolour},   
    commentstyle=\color{codegreen},
    keywordstyle=\color{blue},
    numberstyle=\tiny\color{codegray},
    stringstyle=\color{codepurple},
    basicstyle=\ttfamily\footnotesize,
    breakatwhitespace=false,         
    breaklines=true,                 
    captionpos=b,                    
    keepspaces=true,                 
    numbers=left,                    
    numbersep=5pt,                  
    showspaces=false,                
    showstringspaces=false,
    showtabs=false,                  
    tabsize=2
}

\lstset{style=mystyle}


\begin{document}

\section{Ler 3 numeros, imprimir em ordem e a soma dos dois menores}

Primeiro comecamos lendo os numeros do usuario e salvando cada um em uma variavel

    \begin{lstlisting}[language=Python]
    n1 = int(input("Digite o primeiro numero")
    n2 = int(input("Digite o segundo numero")
    n2 = int(input("Digite o terceiro numero")\end{lstlisting}

(A funcao \textit{int()} transforma o valor entre parenteses em um numero. Isso e necessario pois a funcao \textit{input()} retorna um valor de tipo texto(string)).

Agora vamos criar mais tres variaveis para armazenar os valores ordenados.

\begin{lstlisting}[language=Python]
    maior = 0   #Iniciamos com 0 pois em python
    meio = 0    #nao podemos iniciar variaveis sem valor
    menor = 0\end{lstlisting}


Agora finalmente vamos ordenar os numeros. Precisaremos de tres conjuntos de \textit{if/else}, uma para caso o maior for n1, uma pro n2 e outr pro n3. Cada um dessas estruturas contendo mais duas, para ordenar os dois numeros restantes


\begin{lstlisting}[language=Python]
    if(n1 > n2 and n1 > n3):
        maior = n1
        if(n2 > n3):
            meio = n2
            menor = n3
        else:
            meio = n3
            menor = n2
    if(n2 > n1 and n2 > n3):
        maior = n2
        if(n1 > n3):
            meio = n1
            menor = n3
        else:
            meio = n3
            menor = n1
    if(n3 > n1 and n3 > n2):
        maior = n3
        if(n1 > n2):
            meio = n1
            menor = n2
        else:
            meio = n2
            menor = n1\end{lstlisting}

Note que poderiamos ordenar diretemente nas variaveis iniciais, utilizando apenas mais uma variavel auxiliar, ou ate mesmo imprimir os os resultados diretamente. Porem, vamos fazer dessa forma para facilitar o entendimento.

Pronto. agora e apenas somar e imprimir 

\begin{lstlisting}[language=Python]
    print(f"Ordem decrescente: {maior}, {meio}, {menor}\n")
    print(f"A soma dos menores e: {menor} + {meio} = {menor+meio}\n")\end{lstlisting}

\end{document}

